
\subsection{算子 $\mathfrak q$ 的构造 I：$A_{n,K}$ 版本}
\label{subsec:7-4-2}
\liftIndex{I387}\liftIndex{I359}

本小节讨论命题~\ref{prop:3-8-66}。
事实上，
为解决第~\ref{subsec:7-2-3} 小节
所述的问题，
我们须先修改原有构造：
在某个有限阶段停止构造，
转而使用同调代数。
我们将用这一修改后的构造
完成定理~\ref{thm:3-8-32} 的证明。
修改后的版本即下面的
引理~\ref{lem:7-4-17}；
其证明与第~\ref{sec:7-2} 节的论证类似。

设 \(\mathcal X_1(L)\) 是起始的
\(L\) 上奇异单形的可数集。
假定对
\(P_i\in\mathcal X_1(L)\)、
\((\beta,k)\lesssim(n,K)\)，
\(\mathcal M^{\mathrm{main}}_{k+1}
(L;\beta;\vec P)\)
上已有一个多值截面系统；
\(\lesssim\) 见定义~\ref{def:7-2-63}。
置
\(\mathcal X_{1,00}(L)=\mathcal X_1(L)\)。
假设此系统定义了一个
滤过 \(A_{n,K}\) 结构，
且满足性质~\ref{prop:7-2-27}。
再设 \(\mathcal X(M)\) 是
\(M\) 上满足假设 3.8.28 的
奇异单形可数集。
为构造命题~\ref{prop:3-8-66} 中的
\(\mathcal X_{1+}(L)\)，
需要找到一列链的可数集
\(\mathcal X_{1+,g}(L)\)，满足
\[
\mathcal X_{1+,0}(L)
\supset\mathcal X_{1+,00}(L)
=\mathcal X_1(L).
\]
记
\[
\mathcal X_{1+,(g_0)}(L)
=\bigcup_{g=0}^{g_0}
\mathcal X_{1+,g}(L).
\]
若 \(P\in\mathcal X_{1+,g}(L)\)，
便称 \(P\) 为第 \(g\) 代。
与第~\ref{sec:7-2} 节相同，
定义带装饰的模空间
\[
\mathcal M^{\mathrm{main},\mathfrak d}_{k+1,\ell}
(L;\beta;\vec Q,\vec P).
\]
注意装饰仅是映射
\(\mathfrak d\colon
\{1,\ldots,k\}\to\mathbb Z_{\ge0}\)：
只有 \(P\) 具有代数，
因此不存在与 \(Q\) 对应的装饰。

\begin{lemma}{7.4.17}\label{lem:7-4-17}
对任意 \(\delta>0\) 与 \(K>0\)，
存在 \(g\le K\) 时的
\(\mathcal X_{1+,g}(L)\)，
以及在
\(\|(\mathfrak d,\beta)\|\le K\)、
\(\ell<K\) 且
\(P_i\in\mathcal X_{1+,(K)}(L)\)
时的多值截面系统
\(\mathfrak s_{\mathfrak d,k,\beta,
\vec Q,\vec P}\)，
使各截面横截于零。
\translatornote{源书此处将
\(\|(\mathfrak d,\beta)\|\)
写作参数反序的
\(\|(\beta,\mathfrak d)\|\)，
而第 7.2 节仅定义前者，
本引理的条件 (7.4.18.2)
也用前者；见疑误 A-0561。
源书还将输入集合写作
\(\mathcal X_{1,(K)}(L)\)，
但本小节只定义了
\(\mathcal X_{1+,(K)}(L)\)；
此处据此补加号，见疑误 A-0560。}
此外，满足：

\par\noindent\textup{(7.4.18.1)}\label{cond:7-4-18-1}\quad
\(\mathcal X_{1+,g}(L)\)
满足性质~\ref{prop:7-2-27}。

\par\noindent\textup{(7.4.18.2)}\label{cond:7-4-18-2}\quad
若
\(P_i\in\mathcal X_{1+,\mathfrak d(i)}(L)\)
（\(i=1,\ldots\)），则
\[
\begin{aligned}
\operatorname{ev}_{0*}\bigl(
&\mathcal M^{\mathrm{main},\mathfrak d}_{k+1,\ell}
(L;\beta;\vec Q,\vec P)^{\mathfrak s_{
\mathfrak d,k,\beta,\vec Q,\vec P}}\bigr)
\end{aligned}
\]
可分解为
\(\mathcal X_{1+,(g)}(L)\) 中的元素，
其中
\(g=\|(\mathfrak d,\beta)\|\)。

\par\noindent\textup{(7.4.18.3)}\label{cond:7-4-18-3}\quad
多值截面
\(\mathfrak s_{\mathfrak d,k,\beta,
\vec Q,\vec P}\)
按与第~\ref{sec:7-2} 节相似的意义相容。

\par\noindent\textup{(7.4.18.4)}\label{cond:7-4-18-4}\quad
若 \(\ell=0\) 且
\(P_i\in\mathcal X_{1+,00}(L)
=\mathcal X_1(L)\)，
则
\(\mathfrak s_{\mathfrak d,k,\beta,
\vec Q,\vec P}\)
与起始的多值截面系统一致。

\par\noindent\textup{(7.4.18.5)}\label{cond:7-4-18-5}\quad
\[
\operatorname{ev}_{0*}\bigl(
\mathcal M^{\mathrm{main},\mathfrak d}_{k+1,\ell}
(L;\beta;\vec Q,\vec P)^{\mathfrak s_{
\mathfrak d,k,\beta,\vec Q,\vec P}}\bigr)
\]
位于原模空间
\(
\mathcal M^{\mathrm{main},\mathfrak d}_{k+1,\ell}
(L;\beta;\vec Q,\vec P)
\)
的 \(\delta\)-邻域内。
\translatornote{源书将
\(\operatorname{ev}_{0*}\) 推前后的链
与未推前的模空间直接比较邻域，
二者所在空间不同；
此处保留原式，具体应如何定型待核。
见疑误 A-0562。}
\end{lemma}

引理~\ref{lem:7-4-17}
的证明使用第~\ref{sec:7-2} 节
已经多次讨论过的同一归纳过程，
细节从略。

引理~\ref{lem:7-4-17}
给出了定理~\ref{thm:3-8-62}
的一个“\(A_{n,K}\) 版本”
（后文将以引理 7.4.113 精确陈述）。
为了从它证明 \(A_\infty\) 版本，
需要用到与第~\ref{sec:7-2} 节
相当不同的同调代数方法。
这些方法将在第 7.4.3--7.4.4 小节讨论，
定理~\ref{thm:3-8-62}
的证明将在第 7.4.5 小节完成。

导言中的定理~\ref{thm:intro-Y}
与定理~\ref{thm:3-8-62} 相关，
但内容更繁复；
我们将在第 7.4.6 小节证明它。
第 7.4.5 小节采用“配边论证”，
第 7.4.6 小节采用“分叉论证”
（关于这两种方法，见
第~\ref{subsec:7-2-14} 小节）。
第 7.4.6 小节的结果蕴含
第 7.4.5 小节的结果；
我们仍分别叙述，
以表明两种方法均可用于证明
定理~\ref{thm:3-8-62}。

\begin{remark}{7.4.19}\label{rem:7-4-19}
我们感谢 Cheol-Hyun Cho 指出：
证明定理~\ref{thm:3-8-62}
需要使用与第~\ref{sec:7-2} 节
所用方法相当不同的同调代数。
\end{remark}
